% Supplementary material: moved from Sections 6.1-6.2 of the main text.

\section{Platforms and Device Tiers in Detail}
\label{supp:deployment}

\subsection{Apple Intelligence (AFM)}

Apple's Foundation Models \citep{apple2024intelligence,gunter2025afm} are the most fully documented on-device \gls{slm} deployment. The 2025 report \citep{gunter2025afm} describes a $\sim$3B on-device model, optimized for Apple silicon through KV-cache sharing (dividing the model so the second block reuses the first block's cache, a reported 37.5\% memory saving) and 2-bit quantization-aware training, paired with a Parallel-Track Mixture-of-Experts server model on Private Cloud Compute. The on-device model is reported competitive with Qwen3-4B and Gemma-3-4B on its target tasks, and the server model with Llama~4 Scout (\evC). The third-generation family \citep{apple2026afm3}, five models built in collaboration with Google, adds a 20B-total on-device model that activates 1--4B parameters per token and keeps its full weights in flash, swapping experts into DRAM as needed; Private Cloud Compute was extended to NVIDIA GPUs in Google Cloud. This followed a January 2026 joint statement that Apple's next foundation models would be based on Gemini models and cloud technology \citep{google_apple2026statement}. Google's Gemini Nano, the on-device member of the Gemini family \citep{gemini2023}, occupies the same design point on Android.

\subsection{Samsung Gauss}

Samsung announced Gauss in November 2023 as a family of generative models in three parts, Gauss Language, Gauss Code, and Gauss Image, used first for internal productivity and planned for its products \citep{samsung2023gauss}. Samsung has published no parameter count, architecture, or on-device specification for it, so Gauss fails criterion M1, is not in the inventory, and we make no claim about its deployment. It is listed here because it is frequently cited as an on-device \gls{slm}, and the published record does not support that description.

\subsection{General Mobile Deployment Considerations}

\glspl{slm} targeting mobile deployment must navigate three constraints. (1) \emph{Memory}: Consumer smartphones typically provide 6--12~GB RAM, of which the OS uses 2--4~GB, leaving 4--8~GB for applications. A 3B model at Q4 requires ${\sim}2$~GB, comfortable for most flagship devices. (2)~\emph{Thermal}: Sustained inference on mobile \glspl{npu} is throttled by thermals; realistic sustained throughput is lower than peak throughput benchmarks suggest. (3)~\emph{Battery}: published per-device battery figures are scarce and not reported under a common protocol. The fundamental constraint in edge decoding is memory bandwidth rather than compute: flagship phones and mainstream laptops offer on the order of 50--100~GB/s (single-board computers well under 30~GB/s, unified-memory laptop designs considerably more) versus 2--3~TB/s for datacenter GPUs \citep{xu2024device}, making INT4 quantization not merely a storage optimization but a performance enabler for single-stream generation.

% ----------------------------------------------------------------
\subsection{Edge and IoT Deployment}
\label{subsec:deployment_edge}

Below smartphone-class hardware lies a diverse ecosystem of edge devices: Raspberry Pi single-board computers, NVIDIA Jetson modules, and microcontrollers. \glspl{slm} targeting this tier typically require:

\begin{description}
    \item[Sub-1B or extreme quantization] A Raspberry Pi~5 (8~GB RAM, Cortex-A76 CPU, about 17~GB/s of memory bandwidth) can hold 4-bit models up to a few billion parameters, and its decode rate is bounded by that bandwidth. A benchmark of runtimes on single-board computers reports that llamafile reaches up to $4\times$ the throughput of Ollama at 30--40\% lower power for equivalent models \citep{nguyen2025sbc} (\evB).

    \item[Jetson Orin] NVIDIA's Jetson Orin modules, with TensorRT-LLM \citep{nvidia_tensorrt2024} among their runtimes, hold quantized 3--7B models in 8--64~GB of unified memory, with the larger-memory NX and AGX parts suited to the upper end of that range, for industrial and robotics use.

    \item[Microcontrollers (MCU)] ARM Cortex-M class microcontrollers carry roughly 512~KB--2~MB of SRAM. Even a 10M-parameter model needs about 5~MB at 4 bits, and a 100M-parameter one about 50~MB, so a language model of the size this survey covers cannot be resident in microcontroller SRAM, and would at best stream weights from external memory at a large cost in speed. Keyword detection and sensor analytics on these parts are served by task-specific networks far smaller than any model in the inventory; we found no documented \gls{slm} deployment at this tier, and it is outside the scope of this paper.
\end{description}

Table~\ref{tab:appendix_hardware} details the specifications of the four
representative devices referenced throughout this section, from the Raspberry~Pi
5 up to a laptop-class CPU.

% ----------------------------------------------------------------
